% Source: Confluence page 14647353, version 2.
\chapter{Introduction}
\label{ch:introduction}
\section{Background and Motivation}

Planning a short trip or local outing is a common activity, whether for leisure, social gatherings, sightseeing, or exploring new places. Unlike extended vacations, these activities typically take place within a few hours or a single day. Although their duration is limited, organizing a suitable itinerary can still require considerable effort, particularly when travelers must account for available time, budget, transportation, personal preferences, and the locations they wish to visit.


Digital mapping and travel-planning applications have made it easier to discover destinations, obtain directions, and organize travel information. However, planning a complete outing often involves combining information from multiple sources and making decisions about which activities to include, how to arrange them, and whether the resulting schedule is realistic. These decisions become more difficult when users have limited time or have not yet decided where they want to go.


Furthermore, an itinerary that appears suitable before departure may become impractical during the trip. Unexpected delays, changes in activity duration, or decisions to skip or add destinations can affect the remaining schedule. Adjusting a plan under these circumstances may require users to reconsider their available time, travel routes, and priorities.


These challenges motivate the development of Trippa, a smart short-trip planning system intended to simplify itinerary creation and management. The project focuses on providing an accessible planning experience that considers practical constraints while allowing users to adjust their plans as circumstances change.


\section{Problem Statement}

The main problem addressed by this project is the difficulty of creating and maintaining a realistic itinerary for a short trip with limited planning effort.


Users may know how much time and money they have available without knowing which destinations would best suit their circumstances. Manually identifying suitable places, estimating travel durations, arranging activities, and checking schedule compatibility can be time-consuming. Moreover, selecting individually attractive destinations does not necessarily result in a practical itinerary when travel time, opening hours, and activity durations are considered together.


Another challenge is the difference between a planned itinerary and the actual trip experience. Changes occurring during an outing can invalidate the original schedule, leaving users to determine which remaining activities are still possible.


The project therefore addresses two related needs: generating an itinerary that satisfies the user's initial constraints and supporting adjustments when the trip no longer follows the original plan.


\section{Project Objectives}

The primary objective of Trippa is to develop a short-trip planning system that helps users create, manage, and follow practical itineraries with minimal manual effort.


The project aims to achieve the following objectives:


\begin{enumerate}
\item Simplify trip planning: Allow users to create itineraries using basic information such as available time, location, budget, transportation, and personal preferences.
\item Generate feasible itineraries: Produce organized schedules that consider travel durations, activity durations, opening hours, and other relevant constraints.
\item Support flexible customization: Enable users to review and modify generated itineraries according to their needs.
\item Assist during trip execution: Provide a way to follow trip progress and adjust remaining activities when unexpected changes occur.
\item Improve accessibility and usability: Offer an intuitive interface that supports convenient planning and trip management across desktop and mobile devices.
\end{enumerate}


Together, these objectives seek to reduce planning complexity while improving the practicality and flexibility of short-trip itineraries.


\section{Project Scope}

Trippa primarily targets local outings and single-day trips within an urban area. The intended users include individuals or small groups planning recreational activities, casual outings, or short visits to multiple destinations.


The system covers the main stages of a short trip, from collecting user requirements and generating an initial itinerary to reviewing, modifying, and following the plan. It also considers changes during an ongoing trip and supports adjustments to the remaining schedule. Mapping and location information are incorporated to improve the practicality of generated itineraries, while artificial intelligence may provide additional assistance in interpreting user requests and interacting with the planning process.


The project emphasizes itinerary planning and management rather than comprehensive travel services. Functions such as flight and accommodation booking, payment processing, complex multi-day vacation management, and dedicated turn-by-turn navigation are outside its initial scope.


Detailed functional requirements, system boundaries, and technical considerations are presented in subsequent chapters.


\section{Expected Contributions}

The expected outcome of this project is a functional software system that demonstrates how multiple planning considerations can be integrated into a unified short-trip planning experience.


From a user's perspective, Trippa is intended to reduce the effort required to organize an outing, provide clearer information about the planned schedule, and offer greater flexibility when plans change. Rather than requiring users to manually coordinate multiple destinations, the system aims to assist them in making practical decisions based on their available resources and preferences.


From a software engineering perspective, the project provides an opportunity to apply requirements analysis, system modeling, full-stack application development, geographic service integration, and constraint-based itinerary planning to a practical problem. It also explores how optional AI-assisted interactions can complement a structured planning system.
